\chapter{结论与展望}
\label{chap:conclusion}

本文以 3/2 随机波动率模型下的欧式期权定价问题为核心研究对象，针对现有研究中定价方法性能评价多局限于理想化学术基准算例、缺乏真实极端市场环境实证检验的不足，系统构建了 3/2 模型下十一类欧式期权定价方法的完整理论框架，通过多场景数值实验量化了不同方法的定价精度、计算效率与数值稳定性，并基于 2025 年 10 月 BTC 市场闪崩极端事件，完成了模型的全周期参数校准与定价方法的极端场景再检验，厘清了不同方法在理想化环境与真实市场中的性能边界，为 3/2 模型在衍生品定价中的理论应用与工程实践提供了系统性的分析框架与实证依据。本章将系统总结全文的核心研究结论，客观分析研究存在的局限性，并对未来相关研究方向进行展望。

\section{主要研究结论}

本文通过理论构建、数值实验与实证分析三个层面的系统研究，形成了四项核心研究结论，完整覆盖了 3/2 模型的定价方法性能、市场适配性与实践应用边界。

第一，本文系统厘清了纯扩散框架下 3/2 模型十一类期权定价方法的理论逻辑与性能边界，形成了清晰的方法性能分层体系。本文将十一类定价方法划分为时间步进类条件蒙特卡洛方法、终端精确与近似精确模拟方法、傅里叶频域定价方法三大类别，明确了各类方法的实现路径与理论假设，并基于 7 组覆盖常规市场与极端参数环境的学术基准算例，完成了方法的系统性性能评价。研究结果显示，十一类方法可按工程实用性划分为五个梯队：第一梯队为傅里叶频域的 COS 与 FFT 方法，是唯一在全场景保持数值稳定、同时具备极高定价精度与计算效率的方法，其中 COS 方法的综合性能全面领先；第二梯队为时间步进类的精确步进、Poisson-Gamma 与 QE 近似精确步进方法，是蒙特卡洛框架下唯一全场景无数值失效的方案，兼顾了定价精度与计算效率，是需要完整路径生成场景的唯一可靠选择；第三梯队为 Euler/Milstein 离散格式，仅在极细时间步长下能达到可接受的定价精度，计算效率极低，仅可作为理论精度对比的基准工具；第四梯队为 IG 与 Gamma 近似精确模拟方法，仅在特定平值场景下偏差可控，绝大多数场景出现系统性大幅定价偏差，仅可用于特定场景的学术对比；完全失效梯队为 Baldeaux 精确模拟、Fourier-Cosine 终端模拟与矩匹配终端模拟方法，仅能在单一基准场景下输出有效结果，其余场景完全数值失效。核心发现是，定价方法理论上的无偏性与“精确”标称，不等同于工程实践中的全场景稳定性，多数标称“精确”的终端模拟方法仅能在理想化基准参数下完成理论验证，无法适配非基准的极端参数环境。

第二，3/2 随机波动率模型能够有效捕捉加密货币高波动市场的期权定价特征，在极端行情下具备良好的市场适配性。本文构建了“长期参数时间序列估计 - 时变参数横截面校准”的两步校准框架，基于崩盘事件前半年的 DVOL 日度数据完成了长期均值回复参数的估计，基于 Medvedev-Scaillet 渐近展开方法完成了崩盘事件全周期的时变参数校准。实证结果表明，纯扩散框架下的 3/2 模型能够精准拟合崩盘事件全周期波动率微笑的形态演化，模型核心参数呈现出显著的事件驱动型时变规律：崩盘期平值隐含波动率与波动率的波动率同步大幅跳升，6 天短期期限的涨幅显著高于 20 天中期期限，反映出市场对极端事件的短期波动预期与尾部风险定价显著提升；价格与波动率的负相关系数在崩盘期出现阶段性弱化，短期期权合约的杠杆效应异化更为显著，而中期合约的相关关系保持了相对稳定，体现了极端行情下恐慌性交易与流动性冲击对短期期权定价结构的显著影响。这一结果验证了 3/2 模型在加密货币高波动市场中的良好适配性，为加密衍生品市场的定价与风险测度提供了可靠的模型基础。

第三，真实极端市场场景对定价方法的数值稳定性与鲁棒性提出了远高于学术基准算例的要求，进一步放大了不同方法的性能分化。本文基于校准得到的 6 组真实市场参数，构建了全新的极端场景测试案例，完成了定价方法的性能再检验，并与学术基准算例结果进行了横向对比。研究发现，学术基准算例中仅部分场景失效的终端精确模拟方法，在真实极端场景中出现了全场景的数值失效或系统性大幅定价偏差，工程实用性完全丧失；而时间步进类的分布抽样方法与傅里叶频域方法，在学术案例与真实极端场景中均保持了 100\% 的数值稳定性与稳定的定价精度，展现出极强的极端行情鲁棒性。同时，极端场景进一步凸显了 COS 方法的性能优势，其在全场景保持了趋近于 0 的定价偏差，单案例计算时间稳定在 0.05 秒左右，是极端市场环境下欧式期权批量定价、实时估值与模型校准的最优方案；Poisson-Gamma 与 QE 近似精确步进方法则是路径依赖期权定价与动态对冲场景的最优选择，彻底弥补了传统 Euler/Milstein 格式步长敏感、极端场景精度极差的核心缺陷。

第四，本文基于理论与实证结果，形成了 3/2 模型期权定价方法的场景化应用指引。对于欧式期权批量定价、模型参数校准等场景，优先选择 COS 频域方法，其全场景稳定性、最高精度与最优效率可充分适配此类场景的需求；对于路径依赖期权定价、动态对冲等需要完整方差路径生成的场景，优先选择 Poisson-Gamma 或 QE 近似精确步进方法，其全场景稳定性与精度效率平衡可满足此类场景的工程需求；对于理论精度基准对比场景，仅可在基准参数环境下使用 Baldeaux 精确模拟方法；其余终端类方法不推荐在实际定价场景中使用，仅可用于特定场景的学术对比分析。

\section{研究局限性}

本文围绕 3/2 随机波动率模型的期权定价问题开展了系统性的理论与实证研究，形成了较为完整的分析框架，但受研究范围、数据条件与模型设定的限制，仍存在一定的局限性，主要体现在以下四个方面：

第一，模型设定的维度存在拓展空间。本文聚焦于纯扩散框架下的 3/2 随机波动率模型，未在模型中引入标的价格跳跃项、随机利率与随机流动性等市场摩擦因素。尽管纯扩散模型已能较好地拟合崩盘事件全周期的波动率微笑形态，但极端行情下标的价格的跳跃特征、市场流动性的瞬时变化，仍可能对期权定价产生系统性影响，模型设定对真实市场结构的刻画仍有进一步丰富的空间。

第二，实证研究的覆盖范围有待拓宽。本文的实证分析仅基于 Deribit 交易所的 BTC 期权数据，未覆盖 ETH 等其他主流加密资产，也未拓展至传统金融市场的股票、大宗商品、外汇等标的。不同标的资产的波动特征、市场微观结构与投资者结构存在显著差异，3/2 模型在不同市场中的适配性与参数特征，仍需更多维度的实证数据进行交叉验证。

第三，定价方法的评价范围存在局限。本文的定价方法性能评价仅聚焦于欧式期权，未针对美式期权、障碍期权、亚式期权等主流路径依赖型衍生品开展分析。3/2 模型在路径依赖类期权定价中的方法构建、精度优化与性能评价，是现有研究未充分覆盖的内容，也是本文研究范围的核心局限所在。

第四，参数校准框架的丰富度仍有提升空间。本文采用的 Medvedev-Scaillet 渐近展开方法更适用于短期到期期权，对于超长期到期期权的校准精度仍有优化空间；同时本文仅采用单一校准方法完成参数估计，未引入基于特征函数的全局校准、贝叶斯校准等方法对参数估计结果进行稳健性对比，校准框架的多元性与稳健性仍有进一步提升的可能。

\section{未来研究展望}

基于本文的研究结论与存在的局限性，未来可从以下五个方向开展进一步的深入研究，完善 3/2 随机波动率模型的定价理论与实践应用体系：

第一，拓展 3/2 随机波动率模型的理论设定。在纯扩散模型的基础上，引入同步跳跃、独立跳跃、随机利率与随机流动性等因素，构建更贴合真实市场结构的扩展 3/2 模型，推导对应模型的特征函数闭式解与期权定价公式，量化不同扩展因素对期权定价的边际影响，进一步提升模型对极端市场环境的适配性；同时开展扩展模型与基准模型的定价精度横向对比，明确不同模型设定的适用场景。

第二，拓展定价方法的适用范围与计算性能优化。针对美式期权、障碍期权、亚式期权等主流路径依赖型衍生品，构建适配 3/2 模型方差过程特性的高效定价方法，完成不同方法的定价精度、计算效率与数值稳定性的系统性评价，形成覆盖全品类衍生品的 3/2 模型定价体系；同时结合 GPU 并行计算、向量化优化等技术，进一步提升大规模蒙特卡洛模拟与频域方法的计算效率，满足实际业务中高频实时定价的需求。

第三，拓展实证研究的市场覆盖维度。将实证分析从 BTC 期权拓展至 ETH 等其他主流加密资产，以及全球股票市场、大宗商品市场、外汇市场等传统金融领域，验证 3/2 模型在不同波动特征、不同市场结构下的适配性，明确模型的通用适用边界；同时基于更长周期的历史数据，开展 3/2 模型与 Heston 模型、SABR 模型等主流随机波动率模型的大样本横向对比，量化不同模型的定价精度与市场解释力差异，为不同市场的模型选择提供实证依据。

第四，优化 3/2 模型的参数校准框架。针对渐近展开方法在长期期权校准中的不足，引入基于特征函数的全局校准、贝叶斯校准、机器学习辅助校准等方法，提升极端市场场景下的参数校准精度与稳定性；同时构建模型参数的实时动态校准体系，实现极端行情下模型参数的高频更新，适配衍生品市场实时风险管理的需求。

第五，基于 3/2 模型构建衍生品全流程风险管理体系。在定价研究的基础上，推导 3/2 模型下期权希腊字母的闭式解与高效数值计算方法，构建基于模型的 VaR、CVaR 等尾部风险测度体系，量化极端行情下期权组合的风险暴露；同时基于模型参数的时变特征，构建波动率风险溢价的测度模型，为加密衍生品市场的对冲策略设计、动态资产配置提供理论支撑与实证依据。